\documentclass[conference]{IEEEtran}
\usepackage{cite}
\usepackage{amsmath,amssymb,amsfonts}
\usepackage{algorithmic}
\usepackage{graphicx}
\usepackage{textcomp}
\usepackage{xcolor}
\usepackage{listings}
\usepackage{booktabs}

\begin{document}

\title{Shift-Left Timing Closure: Pre-Synthesis Structural Diagnostic Engine for Fast RTL Feedback}

\author{\IEEEauthorblockN{Dev Tyagi}
\IEEEauthorblockA{\textit{Independent Researcher} \\
dev@example.com}
}

\maketitle

\begin{abstract}
Traditional Field Programmable Gate Array (FPGA) and Application-Specific Integrated Circuit (ASIC) design flows suffer from a massive feedback latency. RTL engineers rely on post-synthesis Static Timing Analysis (STA) to identify critical paths and logic bottlenecks, a process that can take hours for complex designs. This paper introduces ConstraintForge, an open-source, pre-synthesis structural diagnostic engine that parses Verilog/SystemVerilog into a technology-agnostic Abstract Syntax Tree (AST). By executing a topological Breadth-First Search (BFS) over the generated logic network, ConstraintForge accurately estimates combinational logic depth and high-fanout nets in milliseconds. Experimental results on an Out-of-Order RISC-V core demonstrate that the tool successfully flags critical setup time violations—such as a 28-gate deep instruction fetch unit—in under 0.1 seconds, achieving a 1000x speedup in feedback iteration compared to traditional logic synthesis. 
\end{abstract}

\begin{IEEEkeywords}
Electronic Design Automation (EDA), Static Timing Analysis (STA), Field Programmable Gate Arrays (FPGA), Logic Synthesis, Shift-Left
\end{IEEEkeywords}

\section{Introduction}
The cornerstone of modern digital logic design is Register-Transfer Level (RTL) verification and timing closure. However, the standard methodology presents a "Catch-22" for hardware engineers: one cannot definitively know if a design meets timing requirements without running a full technology mapping and place-and-route (P\&R) pass. Tools such as Xilinx Vivado or Synopsys PrimeTime are highly accurate but computationally expensive. A single compilation pass for a moderately complex System-on-Chip (SoC) can take upwards of several hours. If a designer accidentally writes a deeply nested combinational chain or an excessively wide multiplexer, they will not discover the resulting setup time violation until the end of this compilation cycle.

To address this massive feedback latency, we propose a "shift-left" methodology in hardware design. We introduce \textit{ConstraintForge}, a lightweight command-line interface (CLI) tool designed to run locally in the developer's environment. Rather than waiting for full synthesis, ConstraintForge leverages Yosys as a headless parser to convert raw RTL into a technology-agnostic directed acyclic graph (DAG). It then performs heuristic topological traversal to locate structural bottlenecks before physical constraints are even applied.

\section{Proposed Architecture}
The ConstraintForge diagnostic engine operates in three distinct phases: Graph Extraction, Topological Traversal, and Diagnostic Reporting.

\subsection{Headless AST Extraction}
ConstraintForge invokes the open-source synthesis suite, Yosys, as a headless subprocess. Instead of running a full technology mapping targeting a specific vendor's LUT architecture, the tool halts immediately after the \texttt{proc} and \texttt{flatten} passes. This produces a generic JSON netlist consisting entirely of foundational logic primitives (AND, OR, MUX) and registers (D-Flip-Flops).

\subsection{Topological Traversal Algorithm}
The JSON netlist is loaded into a Python-based diagnostic engine. The engine builds an adjacency list mapping every logical net to its driving cell and sinking endpoints. To detect setup-time violations, the tool identifies all D-Flip-Flop (DFF) Q-pins and primary input ports as source nodes. It then executes a Depth-First Search (DFS) with memoization to find the longest combinational path terminating at a DFF D-pin. 
Paths exceeding a user-defined threshold (e.g., 20 levels of logic) are immediately flagged as setup violations. Similarly, by mapping the sink count of every register, the tool flags high-fanout nets that will inevitably cause routing congestion and hold-time violations.

\section{Experimental Results}
ConstraintForge was evaluated against an Out-of-Order RISC-V microprocessor architecture. 

\subsection{Execution Speed}
Analyzing the \texttt{fetch\_unit.v} module of the RISC-V core using standard synthesis took approximately 4.2 minutes to generate a timing report. In contrast, ConstraintForge extracted the AST and executed the traversal in $0.12$ seconds, a substantial speedup that enables real-time linting within an Integrated Development Environment (IDE).

\subsection{Diagnostic Accuracy}
ConstraintForge successfully identified a critical path originating from the Program Counter (PC) register and terminating at the Instruction Queue interface, calculating a pure logic depth of 28 gates. This exact path was later verified by standard STA to be the primary cause of Worst Negative Slack (WNS) when targeting a 100MHz clock domain. Furthermore, the tool successfully identified a global enable register driving 66 endpoints, correctly diagnosing it as a high-fanout risk requiring register replication.

\section{Conclusion}
We have presented ConstraintForge, a structural diagnostic tool that drastically reduces the feedback loop for RTL engineers. By estimating logic depth and fanout pre-synthesis, developers can resolve architectural bottlenecks in milliseconds rather than hours. Future work will expand the tool's capabilities to include automated clock domain crossing (CDC) analysis and asynchronous reset tree validation.

\begin{thebibliography}{00}
\bibitem{b1} C. Wolf, ``Yosys Open SYnthesis Suite,'' \textit{http://yosyshq.net/yosys/}, 2024.
\bibitem{b2} J. Bhasker and R. Chadha, \textit{Static Timing Analysis for Nanometer Designs}. Springer Science \& Business Media, 2009.
\bibitem{b3} S. Kilts, \textit{Advanced FPGA Design: Architecture, Implementation, and Optimization}. John Wiley \& Sons, 2007.
\end{thebibliography}

\end{document}
